\documentclass[letterpaper,10pt]{article}
\usepackage[empty]{fullpage}
\usepackage{titlesec}
\usepackage{enumitem}
\usepackage[hidelinks]{hyperref}
\usepackage[english]{babel}
\input{glyphtounicode}
\pdfgentounicode=1
\pagestyle{empty}
\addtolength{\oddsidemargin}{-0.55in}
\addtolength{\evensidemargin}{-0.55in}
\addtolength{\textwidth}{1.10in}
\addtolength{\topmargin}{-0.65in}
\addtolength{\textheight}{1.25in}
\raggedright
\setlength{\parindent}{0pt}
\urlstyle{same}
\titleformat{\section}{\vspace{-4pt}\scshape\raggedright\large}{}{0em}{}[\titlerule\vspace{-5pt}]
\newcommand{\entry}[4]{\textbf{#1}\hfill #2\\\textit{#3}\hfill\textit{#4}\vspace{-3pt}}
\newcommand{\items}{\begin{itemize}[leftmargin=0.17in,itemsep=0pt,topsep=2pt]}
\newcommand{\enditems}{\end{itemize}\vspace{-4pt}}
\begin{document}
{\Large\textbf{Giacomo Cappelletto}}\\[-1pt]
Boston, MA \;|\; \href{mailto:giacomo.cappelletto@icloud.com}{giacomo.cappelletto@icloud.com} \;|\; +1 (857) 753-0133 \;|\; \href{https://www.linkedin.com/in/giacomo-cappelletto/}{LinkedIn} \;|\; \href{https://github.com/JJCAPPE}{GitHub} \;|\; \href{https://cv-nu-sage.vercel.app/}{Portfolio}
\section{Education}
\entry{Boston University}{Boston, MA}{B.S. Computer Engineering, GPA: 3.97/4.00}{Expected May 2028}
\section{Experience}
\entry{Banca Mediolanum}{Milan, Italy}{Software / AI Intern, Customer Knowledge}{Jun 2026 -- Aug 2026}
\items
\item Led development of an internal AI-agent platform on Databricks across architecture, full-stack implementation, evaluation, deployment, and stakeholder iteration.
\item Built a Model Serving agent using MLflow ResponsesAgent and LangGraph with governed dynamic SQL, citation-backed document RAG, policy-controlled tools, request-scoped context, evidence review, and MLflow tracing.
\item Built the React/Vite + FastAPI application and prompt-evaluation platform with SSE streaming, authenticated sessions, PostgreSQL/Lakebase, versioned Delta prompts, LLM judging, Unity Catalog controls, and Asset Bundle deployment.
\enditems
\entry{Boston University -- Prof. Brian Kulis}{Boston, MA}{Undergraduate Machine Learning Researcher}{Aug 2026 -- Present}
\items
\item Researching metric learning for human-motion retrieval by learning fixed-length embeddings from temporal 2D/3D pose sequences and adapting contextual-similarity objectives to motion representations.
\item Evaluating contextual and contrastive objectives against retrieval baselines under pose noise, missing joints, temporal jitter, occlusion, and viewpoint shift using Recall@K and mAP.
\enditems
\entry{Boston University College of Engineering}{Boston, MA}{Undergraduate Machine Learning Researcher}{Jan 2026 -- May 2026}
\items
\item Built a Python biomechanics pipeline using MMPose and MotionBERT to stabilize rowing video, infer 2D keypoints, reconstruct 3D skeletons, and quantify per-frame and stroke-averaged kinematics.
\item Trained sequence-to-sequence regression models mapping time-aligned 3D kinematics to rowing force curves and validated predictions against instrumented on-water and erg telemetry.
\enditems
\entry{Societ\`a Cappelletto S.r.l.}{Treviso, Italy}{Software Engineer}{May 2025 -- Oct 2025}
\items
\item Rebuilt an Electron inventory application in Tauri, React, and Rust, reducing app size by 70\%, memory usage by 80\%, startup time by 60\%, and Shopify inventory-search latency by 43\% through parallel GraphQL processing.
\item Shipped SKU search, Firebase-backed audit history, real-time inventory reconciliation, multi-location updates, release distribution, and automatic application updates.
\enditems
\entry{TickIT GmbH}{Remote}{Frontend \& Backend Engineer}{Nov 2024 -- Feb 2025}
\items
\item Shipped Ruby and Next.js/TypeScript features as one of four developers on a 150,000-line event-ticketing platform, including QR ticket generation/validation, purchasing flows, event pages, and attendee access controls.
\item Built customer-spending analytics, forecasting pipelines, and organizer dashboards for operational metrics, inventory forecasting, and event cost planning.
\enditems
\entry{Boston University Men's Rowing}{Boston, MA}{Software Engineer}{Jan 2026 -- May 2026}
\items
\item Built a mobile-first training platform for 50 athletes and coaches with Supabase/Prisma data models, role-based access, proof-validated logging, reporting, leaderboards, and automated weekly recap emails.
\enditems
\section{Selected Project}
\textbf{Deskinator -- Autonomous Desk-Cleaning Robot} \hfill Python, Raspberry Pi\\[-2pt]
\items
\item Built autonomous coverage software using APDS9960 edge sensing, stepper odometry, RANSAC rectangle fitting, boustrophedon path planning, telemetry, and simulation-based performance analysis.
\enditems
\section{Leadership \& Athletics}
\textbf{Boston University Men's Rowing -- NCAA Division I Student-Athlete} \hfill Sep 2024 -- Present\\[-2pt]
Freshman Student-Athlete of the Year (2025); Most Improved Oarsman (2026).
\section{Skills}
\textbf{Languages:} Python, TypeScript/JavaScript, SQL, Rust, C/C++, Ruby, Java\\
\textbf{Frameworks:} React, Next.js, FastAPI, Node.js, Flask/Dash, Ruby on Rails, Tauri\\
\textbf{ML/Data/Cloud:} Databricks, MLflow, LangGraph, TensorFlow/Keras, MMPose, MotionBERT, PostgreSQL, Supabase, Docker, Google Cloud Run, Git
\end{document}
